\documentclass[11pt,a4paper]{article}
\usepackage[margin=24mm]{geometry}
\usepackage{amsmath,amssymb,booktabs,array,longtable,xcolor,hyperref}
\usepackage[T1]{fontenc}
\hypersetup{colorlinks=true,linkcolor=blue,urlcolor=blue}
\newcommand{\code}[1]{\texttt{#1}}
\newcommand{\high}{\textcolor{red!75!black}{\textbf{HIGH}}}
\newcommand{\medium}{\textcolor{orange!85!black}{\textbf{MEDIUM}}}
\newcommand{\low}{\textcolor{green!45!black}{\textbf{LOW}}}

\title{Technical Review of \code{GNSS\_Processing\_fg}:\\
Pipeline, Double-Difference Management, LAMBDA, and Risk Audit}
\author{Code-level review of the LIGO repository}
\date{August 2026 source snapshot}

\begin{document}
\maketitle

\begin{abstract}
This document describes the GNSS factor-graph pipeline implemented primarily in
\code{src/GNSS\_Processing\_fg.cpp}.  It concentrates on rover/base
double-difference (DD) code and carrier factors, ambiguity-arc management,
wide-lane fixing, raw-ambiguity LAMBDA/PAR, cycle-slip and residual screening,
and fixed-lag factor deletion.  It also records algorithmic defects and risks
found by static inspection.  The most serious issues are: (i) DD observations
sharing a reference satellite are modelled as independent, so ambiguity
covariance and LAMBDA validation can be overconfident; (ii) committed wide-lane
and integer-fix factors are not registered with the fixed-lag ownership lists,
allowing stale ambiguity variables and constraints to accumulate; and (iii) the
wide-lane constraint is committed before the raw fixing stage and has no
measurement-cost rollback test.  These are credible false-fix or long-run
stability risks and should be addressed before treating the FIXED flag as a
safety-critical integrity statement.  Section~\ref{sec:high-remediation}
records the implemented remediation and August 2026 verification of every
HIGH finding; the audit table is retained as the pre-remediation evidence.
\end{abstract}

\tableofcontents

\section{Scope and state convention}

The main graph symbols are:
\begin{center}
\begin{tabular}{lll}
\toprule
Key & Meaning & Mode \\
\midrule
$A(k)$ & local position and velocity & LiDAR enabled \\
$R(k)$ & attitude & both \\
$O(k)$ & angular rate, acceleration, and IMU biases & LiDAR enabled \\
$G(k)$ & gravity & LiDAR enabled \\
$F(k)$ & ECEF position, velocity, accelerometer bias, gyro bias & no LiDAR \\
$B(k)$ & four constellation receiver-clock biases & both \\
$C(k)$ & common receiver-clock drift & both \\
$E(0)$ & ECEF anchor & LiDAR enabled \\
$P(0)$ & local-to-ECEF rotation & LiDAR enabled \\
$n(j)$ & one DD carrier ambiguity, stored in metres & RTK DD \\
\bottomrule
\end{tabular}
\end{center}

The DD sign convention is target minus reference.  For target satellite $s$,
reference $r$, rover $u$, and base $b$, the geometric DD is
\begin{equation}
g_{sr}=\rho_u^s-\rho_b^s-\rho_u^r+\rho_b^r.
\end{equation}
The code factor residual is $g_{sr}-P_{sr}^{DD}$.  The carrier factor residual
is
\begin{equation}
e_{\Phi}=g_{sr}+a_{sr}-\Phi_{sr}^{DD},
\end{equation}
where graph variable $a_{sr}$ is in metres.  LAMBDA converts it to cycles using
$N_{sr}=a_{sr}/\lambda$.

\section{Top-level processing pipeline}

\subsection{Input filtering and initialization}

\code{processGNSS()} passes each rover epoch to
\code{GNSSAssignment::processGNSSBase()}.  That routine:
\begin{enumerate}
  \item accepts GPS, GLONASS, Galileo, and BeiDou only;
  \item requires a primary-frequency observation and a valid ephemeris;
  \item rejects observations exceeding configured pseudorange, Doppler, or
        carrier standard-deviation thresholds;
  \item maintains satellite tracking and Hatch-smoothed pseudorange state;
  \item after initialization, applies an elevation mask and an azimuth-bin
        diversity selection; and
  \item limits the accepted epoch to at most 15 satellites.
\end{enumerate}

Before graph initialization, epochs with fewer than four valid satellites are
discarded.  Valid GNSS epochs and corresponding local pose/velocity samples fill
an initialization window.  \code{GNSSLIAlign()} estimates the ECEF anchor and
local-to-ECEF rotation, then \code{SetInit()} creates priors and the first graph
state.  The initialization has a low-excitation SPP/ENU path and an
excited-motion trajectory-alignment path.

\subsection{Per-epoch optimization}

After initialization, \code{Evaluate()} performs the following sequence:
\begin{enumerate}
  \item construct the current local or ECEF navigation initial estimate;
  \item call \code{AddFactor()} to add rover pseudorange/Doppler, receiver-clock,
        motion, LiDAR/IMU, time-differenced carrier, and optional base/rover DD
        factors;
  \item run a float iSAM2 update through \code{runISAM2opt()};
  \item if carrier DD factors were added, call
        \code{attemptIntegerAmbiguityResolution()};
  \item export the optimized navigation state; and
  \item prune legacy time-differenced carrier history.
\end{enumerate}

This order is conceptually correct: integer validation uses the newly optimized
float state rather than the pre-optimization state.

\section{Double-difference measurement management}

\subsection{Base synchronization and signal candidates}

\code{addDoubleDifferenceFactors()} first finds the closest base epoch within
\code{base\_epoch\_tolerance}.  A missing match skips only the DD branch; the
rover-only graph continues.  For every valid rover satellite and frequency, a
candidate is retained only if:
\begin{itemize}
  \item the frequency maps to a supported raw band (primary, secondary, L5,
        extra, or wide);
  \item carrier phase and carrier standard deviation exist and phase is nonzero;
  \item a base observation matches satellite and frequency within 100 kHz;
  \item neither rover nor base reports an invalid half-cycle; and
  \item the propagated single-difference variance is finite.
\end{itemize}

The single-difference phase and variance are
\begin{align}
\Phi_{u-b}^s &= \lambda_u\phi_u^s-\lambda_b\phi_b^s,\\
\sigma_{u-b,s}^2 &= (\lambda_u\sigma_{\phi,u})^2+
                    (\lambda_b\sigma_{\phi,b})^2.
\end{align}
The carrier DD standard deviation used by each scalar factor is
\begin{equation}
\sigma_{DD,sr}=\sqrt{\sigma_{u-b,s}^2+\sigma_{u-b,r}^2+
                           \sigma_{\mathrm{floor}}^2}.
\end{equation}

\subsection{Cycle-slip screening}

LLI bit 1 (half-cycle ambiguity) rejects a candidate immediately.  LLI bits 0
or 1 set the loss-of-lock flag and cause a new ambiguity arc.  An optional
geometry-free detector forms, for each dual-frequency satellite,
\begin{equation}
L_{GF}=(\Phi_1^u-\Phi_1^b)-(\Phi_2^u-\Phi_2^b)
\end{equation}
and declares a slip when the epoch-to-epoch jump exceeds
\code{rtk\_geometry\_free\_slip\_threshold}.  Continuity also requires the
epoch gap not to exceed \code{rtk\_ambiguity\_gap\_tolerance}.

\subsection{Reference selection}

A reference is selected separately per constellation and raw signal band.  The
lowest single-difference phase variance is preferred, but the previous reference
is retained while available to avoid unnecessary ambiguity redefinition.
Primary and secondary bands are forced to use one common dual-frequency
reference so that $N_{WL}=N_1-N_2$ is well-defined.  This is a good and necessary
constraint.  Other bands can retain independent references.

\subsection{DD pseudorange factors}

Primary-band code DD factors reuse the primary carrier reference.  Their base
Gaussian sigma is \code{double\_difference\_pseudorange\_sigma}.  When
\code{double\_difference\_pseudorange\_robust\_threshold}>0, a Cauchy kernel is
applied in whitened residual units.  Carrier DD factors deliberately remain
Gaussian.  This separation is appropriate because robust carrier weights would
make covariance-based integer validation harder to interpret.

\subsection{Ambiguity arcs and fix-and-hold}

An arc identity is $(r,s,\mathrm{band})$.  A new key $n(j)$ is created for a new
identity, a rover/base LLI event, or a discontinuous epoch gap.  Its initial
value is measured DD minus predicted geometry.  A broad prior uses
\code{ambiguity\_prior\_sigma}.  Each continuous observation increments a lock
counter and stores the latest factor, sigma, elevation, measurement, and
geometry for later screening.

When a reference changes, the code searches a graph of recent fixed DD integer
edges.  If a path connects the new reference and target, it transfers the
equivalent integer and initializes the new arc as held.  The graph-theoretic
sign handling is sensible, but its safety depends completely on old fixed edges
being correct and truly continuous.

\section{LAMBDA and partial ambiguity resolution}

\subsection{Pre-LAMBDA eligibility gates}

Only active, current-epoch, non-fixed ambiguities are eligible.  Each must:
\begin{itemize}
  \item have at least \code{lambda\_min\_lock\_epochs};
  \item have a nonzero wavelength (the DD pair must be integer-compatible);
  \item exist in the current iSAM2 estimate; and
  \item pass the current carrier residual gate.
\end{itemize}
The residual gate evaluates the exact latest factor at the float solution and
rejects the entire target satellite if
$|e_{\Phi}|/\sigma_{DD}>\code{lambda\_max\_normalized\_residual}$.
Satellite-wide rejection is conservative when one frequency is bad.

\subsection{Stage 1: wide-lane fixing}

Primary and secondary ambiguities sharing $(r,s)$ are paired.  A linear
transformation constructs
\begin{equation}
N_{WL}=a_1/\lambda_1-a_2/\lambda_2.
\end{equation}
The full joint marginal covariance is transformed as $Q_{WL}=TQ_aT^T$, which
correctly preserves correlations already represented by iSAM2.  LAMBDA returns
the two best candidates.  The subset must pass a maximum standard-deviation
gate and the ratio test
\begin{equation}
R=\frac{s_2}{\max(s_1,10^{-12})}\geq
\code{lambda\_ratio\_threshold}.
\end{equation}
If it fails, the pair with the largest wide-lane variance is removed until the
PAR minimum is reached.  Accepted integers are imposed by
\code{WideLaneIntegerFactor} and immediately committed to iSAM2.

\subsection{Stage 2: raw ambiguity fixing}

After any wide-lane constraint is committed, the covariance is recomputed;
hence Stage 2 is conditional on the fixed wide lane.  Only primary and secondary
ambiguities are used, grouped by constellation, and only the largest
constellation group is attempted.  L5, extra, and wide raw bands are therefore
float-only in the current Stage 2.

Three PAR deletion orders are attempted independently: lowest elevation,
largest variance, and largest fractional distance from the nearest integer.
For every subset, the code:
\begin{enumerate}
  \item obtains the joint ambiguity covariance;
  \item converts mean and covariance from metres to cycles;
  \item rejects excessive marginal standard deviation;
  \item runs LAMBDA for two candidates and applies the ratio test;
  \item adds proposed integer priors to a copy of iSAM2; and
  \item compares selected current carrier-factor costs before committing.
\end{enumerate}
Accepted raw integer priors implement fix-and-hold.

\subsection{Wide lane versus narrow lane}

The implementation has a real wide-lane stage, but it has \emph{no explicit
narrow-lane observable or narrow-lane integer stage}.  Although
\code{RtkSignalBand::NarrowLane} and a cascade-stage helper exist, no narrow-lane
candidate or factor is constructed in this pipeline.  Stage 2 fixes the original
$N_1$/$N_2$ ambiguities after conditioning on $N_1-N_2$; that is mathematically
related to resolving the short-wavelength information, but it should not be
described in logs or documentation as an implemented narrow-lane combination.

\section{Fixed-lag management}

Each normal factor receives an iSAM2 factor index tracked in
\code{factor\_id\_frame}.  When the active window exceeds
\code{delete\_thred}, factors owned by the oldest frames are removed and
selected boundary priors are inserted.  Legacy time-differenced carrier history
is pruned by frame and time.

Ambiguity state maps and integer-hold factors follow a different lifecycle, and
this difference is the source of one of the high-severity findings below.

\section{Algorithmic flaw and risk audit}

The term ``deadly'' is interpreted here as a flaw capable of producing a false
fixed solution, silent estimator corruption, numerical failure, or unbounded
long-run resource growth.  Static inspection cannot prove field failure rates;
replay testing with truth and fault injection remains necessary.

\begin{longtable}{>{\raggedright\arraybackslash}p{18mm}p{43mm}p{93mm}}
\toprule
Severity & Finding & Evidence, consequence, and recommendation \\
\midrule
\endfirsthead
\toprule
Severity & Finding & Evidence, consequence, and recommendation \\
\midrule
\endhead

\high & Shared-reference DD correlation is omitted &
Every target-minus-reference DD is added as an independent scalar Gaussian
factor (around lines 1754--1761 and 1898--1915).  All DDs in a group reuse the
same reference observation, so their noise covariance has nonzero off-diagonal
terms containing the reference variance.  Adding the reference variance to
each scalar diagonal does not model that correlation.  iSAM2 can consequently
report an overconfident or distorted joint ambiguity covariance, directly
affecting the standard-deviation gate and LAMBDA ratio.  Use a vector DD factor
with the full covariance, whiten a single-difference vector before differencing,
or use an undifferenced formulation with explicit receiver/base clocks. \\
\addlinespace

\high & Integer and wide-lane hold factors escape fixed-lag ownership &
Normal DD factors append their indices to \code{factor\_ids}, but Stage-1
wide-lane factors and Stage-2 integer priors are committed directly via
\code{isam.update()} (around lines 2274--2286 and 2584--2591).  Their returned
indices update \code{id\_accumulate} but are not placed in
\code{factor\_id\_frame}.  The \code{ambiguities\_} and
\code{wide\_lane\_fixes\_} maps are also cleared only on full reset.  Old
ambiguity variables and hold factors can therefore remain indefinitely after
their measurement factors leave the lag window.  This risks unbounded memory
and Bayes-tree growth, and stale state bookkeeping.  Give every ambiguity arc
an ownership/lifetime record and remove its measurement, prior, wide-lane, and
integer-hold factors together after the arc is inactive and outside the window. \\
\addlinespace

\high & Wide-lane fix is irreversible before end-to-end validation &
Stage 1 tests that a trial update is numerically solvable, then commits the
wide-lane constraint immediately.  It does not compare float/fixed measurement
cost, and if Stage 2 later fails the wide-lane factor remains held.  A false
wide-lane ratio pass can thus bias future floats and make later raw fixing look
more precise.  Perform residual/global-cost validation on the trial graph and
commit Stage 1 and Stage 2 atomically, or explicitly retain and roll back the
Stage-1 factor indices when downstream validation fails. \\
\addlinespace

\high & Hatch/cycle-code screening statistics are global, not per satellite &
\code{GNSSAssignment} stores scalar \code{sum\_d} and \code{sum\_d2}, while
\code{processGNSSBase()} resets and updates them inside processing for many
satellites (notably lines 470--497).  The six-sigma test for one satellite can
therefore use statistics last modified by another satellite.  This can cause
incorrect resets or failure to reset a bad track, contaminating smoothed code
and all downstream pseudorange factors.  Store running count, mean/M2, last
phase, and Hatch state per satellite and frequency; use a numerically stable
per-track variance update. \\
\addlinespace

\high & Post-fix validation observes only selected carrier DD factors &
The trial raw fix cost sums only the latest factors of the ambiguities being
fixed (around lines 2528--2558).  It does not check the total graph error,
navigation displacement, pseudorange/DD-code residuals, IMU/LiDAR factors, or
unselected carrier DDs.  The navigation state can absorb a wrong integer while
keeping selected carrier residuals small.  Compare robust global or carefully
partitioned costs, enforce a position jump/innovation bound, and monitor all
current GNSS residuals before commit. \\
\addlinespace

\medium & Local-mode DD lever arm is frozen outside the attitude state &
The caller pre-rotates the antenna lever arm using the input rover rotation;
the local DD factor then uses $p+\ell$ and has no key/Jacobian for $R(k)$
(carrier factor lines 33--56; pseudorange is analogous).  Yet $R(k)$ is an
optimized graph variable.  If attitude changes during optimization, DD geometry
is inconsistent with the optimized antenna position.  Connect local DD factors
to $R(k)$ and use $p+R(k)T_{imu\rightarrow ant}$ with the proper Jacobian, or
formally document and enforce attitude as fixed for these factors. \\
\addlinespace

\medium & No explicit narrow-lane stage &
The enum and cascade helper suggest narrow-lane processing, but the active code
constructs only $N_1-N_2$ wide-lane constraints followed by raw $N_1/N_2$
LAMBDA.  This is not necessarily incorrect, but it is an implementation/documentation
gap and misses explicit narrow-lane screening.  Either remove the unused
narrow-lane terminology or implement and test the intended cascade. \\
\addlinespace

\medium & Raw Stage 2 uses only the largest constellation group &
After eligibility, ambiguities are grouped by the target satellite constellation
and all but the largest group are discarded for that attempt.  This sacrifices
valid information and can starve smaller constellations indefinitely.  Resolve
each independent constellation group separately, or form one correctly
correlated multi-constellation system when the ambiguity definition permits it. \\
\addlinespace

\medium & Elevation used by PAR is target-only &
DD quality depends on both target and reference elevation, but
\code{latest\_elevation} stores only target elevation.  The elevation deletion
method can preserve a DD with a weak reference.  Score by
$\min(elev_s,elev_r)$ or predicted DD variance, and include both elevations in
diagnostics. \\
\addlinespace

\medium & Slip detector is optional and limited to primary/secondary &
When geometry-free detection is disabled, integrity relies on rover/base LLI
and gap checks.  Receivers or RINEX conversions can miss slips, and L5/extra/wide
bands receive no geometry-free consistency check.  Enable the detector after
noise-aware tuning, add Doppler/phase or Melbourne--W\"ubbena confirmation, and
test simultaneous reference/target slips. \\
\addlinespace

\medium & Historical integer transfer has no independent consistency test &
A new reference definition can inherit an integer through a graph of recent
fixed edges and immediately receives fixed-level lock count and prior.  One
incorrect old edge propagates to new arcs.  Require cycle consistency when
multiple paths exist, validate the transferred integer against the current
float residual, and expire edge history explicitly. \\
\addlinespace

\low & Robust DD code noise uses one constant sigma &
DD pseudorange uses a configured constant sigma plus a Cauchy kernel, rather
than propagating four code uncertainties and reference-induced correlation.
The robust kernel limits some outliers but does not repair covariance modelling.
Propagate rover/base code variances and use the same full-covariance strategy as
carrier DD. \\
\bottomrule
\end{longtable}

\section{What is already done well}

Several safeguards are structurally sound:
\begin{itemize}
  \item integer fixing occurs after a float iSAM2 update;
  \item ambiguity values and full joint covariance are consistently converted
        from metres to cycles before LAMBDA;
  \item half-cycle flags are rejected and loss-of-lock/gap events create new
        keys rather than reusing a contaminated ambiguity;
  \item primary and secondary use the same DD reference for wide-lane formation;
  \item residual, lock-duration, precision, ratio, and PAR-size gates are all
        present;
  \item raw integer proposals are evaluated on a copy of iSAM2 before commit;
  \item carrier factors remain Gaussian while robust weighting is isolated to
        DD pseudorange; and
  \item LAMBDA checks covariance finiteness, symmetry, and positive definiteness
        and has a bounded search iteration count.
\end{itemize}

\section{Recommended remediation order}

\begin{enumerate}
  \item Model the full same-reference DD covariance and add a Monte Carlo test
        demonstrating calibrated ambiguity covariance and false-fix rate.
  \item Make wide-lane/raw fixing transactional with global residual and state
        jump validation.
  \item Implement explicit ambiguity-factor ownership and inactive-arc pruning;
        verify bounded graph size on multi-hour replay.
  \item Replace global Hatch screening statistics with per-satellite,
        per-frequency state.
  \item Correct the local lever-arm attitude dependency and validate analytic
        Jacobians numerically.
  \item Add fault-injection tests: target slip, reference slip, simultaneous
        slips, base epoch gap, reference switch, half-cycle transition,
        multipath code outlier, and wrong wide-lane candidate.
  \item Report integrity metrics beyond FIXED/FLOAT: number fixed, subset
        fraction, ratio, worst ambiguity sigma, maximum normalized residual,
        post-fix total-cost change, navigation jump, and ambiguity age.
\end{enumerate}

\section{HIGH-severity remediation and verification}
\label{sec:high-remediation}

Four HIGH findings have complete remedies.  The ambiguity lifecycle crash is
contained by a verified full-graph DD mode, but bounded-memory fixed lag remains
open and requires migration to \code{IncrementalFixedLagSmoother}:
\begin{enumerate}
  \item \textbf{Shared-reference covariance:} because the active graph uses
  scalar DD factors, each same-reference group now uses the conservative
  diagonal majorizer
  $\sigma_i^2=v_i+m v_r+v_f$.  The difference from the exact block covariance
  is $v_r(mI-\mathbf{1}\mathbf{1}^T)\succeq0$; consequently the scalar graph
  cannot be more informative than the correlated block.  Unit test
  \code{SameReferenceScalarApproximationIsConservative} verifies the PSD
  ordering numerically.  DD code sigma is likewise inflated by $\sqrt m$.
  This intentionally trades efficiency for covariance integrity.

  \item \textbf{Ambiguity/hold lifecycle---contained, not fully closed:}
  factor-only deletion is disabled in DD mode, preventing orphan variables and
  the deterministic \code{n4} crash.  A direct
  \code{ISAM2::marginalizeLeaves} attempt was rejected after testing exposed
  that this graph's unconstrained elimination ordering does not guarantee old
  navigation blocks are Bayes-tree leaves; it produced
  \code{Requested to eliminate a key that is not in the factors}.  DD therefore
  currently retains the full graph.  This is numerically stable for the tested
  bag but grows with runtime.  The remaining production fix is to migrate DD
  updates and key timestamps to \code{IncrementalFixedLagSmoother}, which
  creates the required constrained ordering before marginalization.

  \item \textbf{Transactional wide lane:} Stage~1 now creates a pending factor
  graph only.  No wide-lane factor or bookkeeping is committed unless Stage~2
  raw fixing and all post-fix checks pass.  Wide-lane and raw factors are then
  submitted in one iSAM2 update.  Replay logs show
  \code{WIDE-LANE VALIDATED\_PENDING} without leaked \code{FIXED} holds when the
  downstream gates reject the candidate.

  \item \textbf{Per-track Hatch screening:} count, mean, Welford $M_2$, last
  phase/time, and smoothed code are now stored by $(\mathrm{satellite},
  \mathrm{frequency\ index})$.  The outlier test uses deviation from that
  track's mean rather than an absolute cycle--code value.  Unit test
  \code{WelfordStatisticsRemainIndependentPerTrack} checks independence and
  sample variance for two interleaved tracks.

  \item \textbf{End-to-end post-fix validation:} trial validation retains the
  selected-carrier check and additionally compares error over every existing
  graph factor at float and trial-fixed values.  It also rejects a local/ECEF
  navigation-position jump above
  \code{lambda\_max\_position\_jump\_m}.  Unit test
  \code{RequiresGlobalCostAndPositionJumpToPass} covers both rejection gates.
\end{enumerate}

The RTK core suite contains 36 tests after these additions.  The full mapping
target builds successfully, and the supplied bag
\code{2025-03-17-15-11-46.bag} was replayed completely at $2\times$ rate with
DD enabled.  It reached 2910 received epochs and 1639 accumulated DD
observations without the previous epoch-1742195647.998 crash.

\section{Minimum validation matrix}

Before relying on integer fixes, replay truth-labelled data and record false-fix
rate and time-to-fix under the following matrix:
\begin{center}
\begin{tabular}{lll}
\toprule
Dimension & Cases & Required observation \\
\midrule
Baseline & short, medium, long & float/fixed position error and graph size \\
Satellites & 4--6, 7--10, $>10$ & covariance calibration and PAR subset \\
Dynamics & static, urban vehicle, aggressive turns & lever-arm/attitude sensitivity \\
Signal & open sky, multipath, partial blockage & residual and ratio distributions \\
Timing & aligned, boundary tolerance, missing base epochs & arc continuity behavior \\
Slips & target, reference, dual, silent & detection latency and false-fix rate \\
Reference & stable and forced switches & integer-transfer consistency \\
Runtime & at least several hours & bounded variables, factors, and memory \\
\bottomrule
\end{tabular}
\end{center}

\section{Conclusion}

Four HIGH findings identified by this audit now have implemented remedies and
regression/replay evidence.  The lifecycle crash is safely contained, but its
bounded-memory requirement remains HIGH until a timestamped constrained fixed-
lag smoother replaces raw ISAM2.  The conservative covariance majorizer and
atomic/global validation deliberately reduce fixing aggressiveness in exchange
for integrity.  The remaining MEDIUM/LOW findings and the broader fault-
injection matrix should still be completed before making a formal
safety-critical integrity claim.

\end{document}
